\section{第一阶段深讲：内存系统怎样塑造虚拟内存}

操作系统的内存管理不是“把内存分给程序”这么简单。硬件内存系统由 DRAM、cache、TLB、MMU、内存控制器、NUMA 拓扑和 DMA 可见性共同组成。OS 的地址空间、页表、VMA、COW、page cache、伙伴分配器和回收算法，都是在这些硬件条件下形成的。

\subsection{物理内存首先是一张不连续地图}

初学者容易把物理内存想成从 0 到 N 的连续数组。真实机器启动时，固件会告诉内核哪些范围可用，哪些被固件、设备 MMIO、ACPI 表、内核镜像、initrd、显存窗口或保留区域占用。内核不能把保留区当普通页分配。

\begin{lstlisting}
memory map examples:
  usable RAM
  reserved firmware region
  kernel image and initrd
  MMIO aperture for PCIe devices
  ACPI tables and UEFI memory map
  crash kernel reserved area
\end{lstlisting}

因此页分配器的第一步是解析内存地图，建立可用页框集合，而不是假设所有地址都能清零后使用。每个物理页还要记录状态：是否可分配，属于哪个 NUMA node，是否被保留，是否在 LRU，是否 dirty，是否被 DMA pin，是否有硬件错误。

\subsection{页是硬件和 OS 的折中粒度}

页太小，页表和 TLB 压力大；页太大，内部碎片和权限粒度变差。4 KiB 页是许多系统的基础粒度，但大页和巨大页用于降低 TLB miss。OS 在页粒度上做权限、映射、换入换出、COW 和 page cache，是因为硬件 MMU 也按页翻译。

\begin{longtable}{p{0.18\linewidth}p{0.34\linewidth}p{0.34\linewidth}}
\toprule
粒度 & 优点 & 代价 \\
\midrule
小页 & 权限细，碎片小，按需分配灵活 & 页表项多，TLB 覆盖范围小 \\
大页 & TLB 覆盖范围大，page walk 少 & 分配连续内存难，拆分成本高 \\
cache line & 硬件缓存和一致性粒度 & 伪共享、写放大、对齐敏感 \\
块设备扇区/块 & 持久存储 I/O 粒度 & 与页大小不总是一致 \\
\bottomrule
\end{longtable}

这张表说明一个重要原则：OS 很少只有一个粒度。CPU cache 按 cache line 传输，MMU 按页翻译，文件系统按块管理，设备按队列请求传输。性能问题常出现在粒度转换处。

\subsection{虚拟地址先落到 VMA，再落到 PTE}

VMA 描述一段虚拟地址区间的高层语义：权限、是否匿名、是否文件映射、是否共享、增长方向、文件偏移。PTE 描述某一页当前的硬件翻译：是否 present、物理页号、读写执行权限、dirty/accessed 位等。二者不是同一个东西。

\begin{lstlisting}
virtual access:
  find VMA by fault address
  if no VMA: invalid access
  if VMA permissions reject access: protection fault
  inspect or create PTE
  if needed allocate page or read file page
  install PTE and return
\end{lstlisting}

VMA 是“这段地址应该有什么语义”，PTE 是“这页现在如何被硬件翻译”。一个合法 VMA 可以暂时没有 PTE，因为懒分配或文件页尚未读入。一个 PTE 也必须受 VMA 约束，不能凭空给出 VMA 不允许的写或执行权限。

\subsection{页表层级来自地址位数和稀疏性}

若为整个虚拟地址空间建立一个线性页表，会巨大且浪费。多级页表把虚拟地址拆成若干索引，只为实际使用的区域分配中间页表页。地址空间越稀疏，多级页表越有意义。

\begin{lstlisting}
virtual address split:
  top-level index
  middle-level index
  lower-level index
  page-table index
  page offset
\end{lstlisting}

page walk 的成本来自多次内存访问。TLB 缓存最终翻译，是为了避免每次 load/store 都走多级页表。大页减少层级和 TLB 项数量，但要求地址和物理内存按大页对齐，并牺牲一部分细粒度权限。

\subsection{PTE 位是硬件和内核共同维护的协议}

一个 PTE 通常不仅有物理页号，还包含权限和状态位。不同架构位名不同，但角色相似：present 决定是否有当前翻译，writable 决定写入是否允许，user/supervisor 决定用户态能否访问，execute disable 决定能否取指，accessed/dirty 帮助内核判断冷热和写回。

\begin{longtable}{p{0.2\linewidth}p{0.34\linewidth}p{0.34\linewidth}}
\toprule
PTE 信息 & 硬件用法 & 内核用法 \\
\midrule
present & 缺失时触发 page fault & 懒分配、swap、文件缺页 \\
writable & 写保护检查 & COW、只读映射、mprotect \\
user/supervisor & 用户态访问检查 & 内核页隔离、用户地址空间 \\
execute permission & 取指检查 & W\^X、代码页保护 \\
accessed/dirty & 访问和写入记录 & 回收、写回、冷热判断 \\
physical frame & 最终物理页 & 反向映射、引用计数、迁移 \\
\bottomrule
\end{longtable}

PTE 是典型的软硬件共享数据结构。硬件会读取甚至更新某些位，内核会修改权限、清理位、替换物理页。凡是共享，就有一致性问题：修改 PTE 后，TLB 中的旧翻译是否还有效？其他 CPU 是否也缓存了旧结果？这些问题推导出 TLB flush 和 shootdown。

\subsection{缺页处理是一棵决策树}

page fault handler 不能只看“地址不存在”。它要区分访问类型、VMA 类型、PTE 状态、进程状态和内存压力。

\begin{lstlisting}
handle_page_fault(address, access):
  vma = find_vma(address)
  if vma missing: signal segmentation fault
  if access violates vma permissions: signal protection fault
  if anonymous and pte missing: allocate zero page
  if file-backed and pte missing: read page cache or storage
  if write to COW page: copy or reuse page
  if swapped out: read swap entry
  if allocation fails: invoke reclaim or OOM policy
\end{lstlisting}

同样是 fault，结果可能是成功分配、读文件、复制 COW 页、阻塞等待 I/O、触发回收、杀进程或向用户发送信号。把所有 fault 都理解为“程序错了”，会误解虚拟内存。

\subsection{COW 的关键是同时保护父子和旧 TLB}

fork 之后，父子共享物理页。为了让任何一方第一次写入都被内核截获，父子 PTE 都必须改成只读，并标记 COW。若只改子进程，父进程可以继续写共享页。若改了 PTE 但没有处理 TLB，CPU 可能继续使用旧 writable 翻译。

\begin{lstlisting}
fork COW:
  for each writable private page:
    clear writable in parent PTE
    install read-only child PTE
    mark both COW
    increment physical page refcount
    invalidate stale writable TLB translations
\end{lstlisting}

COW 写 fault 的快慢取决于页是否仍共享。如果引用计数为 1，说明只剩当前进程，直接恢复 writable 即可。否则分配新页、复制内容、更新 PTE、减少旧页引用。这里的引用计数必须和页表更新同步，否则可能释放仍被映射的物理页。

\subsection{TLB shootdown 是多核虚拟内存的税}

一个进程可能在多个 CPU 上运行过，或者同一地址空间的多个线程正在不同 CPU 上执行。内核修改某个地址的 PTE 后，要确保所有可能使用旧翻译的 CPU 都失效对应 TLB 项。

\begin{lstlisting}
shootdown timeline:
  CPU0 modifies PTE
  CPU0 records target CPUs using this mm
  CPU0 sends IPI to CPU1 and CPU2
  remote CPUs run interrupt handler and invalidate TLB entry
  remote CPUs acknowledge
  CPU0 may now free or reuse the old page
\end{lstlisting}

这就是为什么大量小范围 \code{mprotect}、\code{munmap} 或频繁映射变化可能在多核机器上很贵。成本包括 IPI、远端中断处理、等待确认和丢失 TLB 热状态。批量 unmap 和延迟释放的很多设计，都是为了减少这笔税。

\subsection{cache 一致性不等于程序没有数据竞争}

硬件 cache coherence 通常保证同一个物理地址的 cache line 在多个 CPU 间遵守一致性协议。但这不等于普通读写自动形成程序想要的同步。两个 CPU 同时更新共享变量，仍可能丢失更新；一个 CPU 写数据再写 flag，另一个 CPU 先看到 flag 后看到旧数据，也可能在弱内存模型上发生。

\begin{lstlisting}
producer:
  data = 42
  release_store(ready, true)

consumer:
  if acquire_load(ready):
    use(data)
\end{lstlisting}

锁、原子操作和内存屏障的目标是建立跨 CPU 可见顺序，不是简单地“让 cache 生效”。OS 内部的 runqueue、引用计数、页状态、设备队列，都需要这种顺序保证。

\subsection{NUMA 让“内存速度”不再单一}

多 socket x86-64 机器中，CPU 访问本地内存和远端内存延迟不同、带宽不同。OS 要知道 CPU 和内存 node 的拓扑，尽量让任务运行在靠近其内存的位置，或者把页迁移到任务常运行的 node。

\begin{longtable}{p{0.2\linewidth}p{0.36\linewidth}p{0.32\linewidth}}
\toprule
策略 & 目的 & 代价 \\
\midrule
first touch & 谁先写页就从谁本地 node 分配 & 初始化线程位置会影响后续性能 \\
CPU affinity & 让任务少迁移，保持缓存和内存局部性 & 负载不均时可能浪费 CPU \\
page migration & 把热页迁到更近 node & 迁移本身消耗带宽和 CPU \\
interleave & 在多个 node 间分散页 & 降低热点但增加部分远端访问 \\
\bottomrule
\end{longtable}

NUMA 问题常被误判为算法慢。同一程序、同一输入，线程绑核和页放置不同，性能可能差很多。OS 调度器和内存管理必须共同考虑这种硬件拓扑。

\subsection{page cache 把内存和存储连接起来}

文件数据读入内存后，不应每次 read 都访问设备。page cache 用内存缓存文件页，普通 read 从缓存复制，普通 write 先写缓存再标 dirty。这样存储 I/O 被合并、重排和批量化，代价是持久化语义变复杂。

\begin{lstlisting}
file read miss:
  allocate folio in page cache
  map file offset to disk block
  submit block I/O
  wait for completion
  mark folio uptodate
  copy to user

file write:
  copy from user to folio
  mark dirty
  writeback later or on fsync
\end{lstlisting}

page cache 也参与内存回收。clean 文件页可以直接丢弃，将来再从文件读；dirty 文件页要先写回；匿名页不能从文件重读，只能保留、换出到 swap 或丢弃进程。内存管理和文件系统因此不是两套独立机制。

\subsection{回收和 OOM 是内存承诺失败后的策略}

当空闲页不足，内核先尝试回收：丢 clean page cache，写回 dirty 页，换出匿名页，收缩 slab 缓存。如果仍无法满足分配，就进入 OOM 策略，选择牺牲某个进程或让分配失败。

\begin{lstlisting}
allocation slow path:
  try per-cpu page lists
  try buddy free lists
  reclaim clean file pages
  write back dirty pages if allowed
  swap anonymous pages if configured
  compact memory for high-order allocation
  invoke OOM policy if still failing
\end{lstlisting}

不同分配上下文允许的动作不同。中断上下文不能睡眠等待磁盘 I/O；持锁路径不能随意进入会回调同一锁的回收；高阶分配可能因为碎片失败，即使总空闲内存不少。内存管理的难点在这些上下文约束。

\subsection{内存系统的最低验收}

第一阶段读者应该能解释：为什么 VMA 合法但 PTE 可以不存在；为什么 COW 要写保护父子双方；为什么改 PTE 后要处理 TLB；为什么 cache 一致性不等于数据竞争安全；为什么 page cache 同时属于文件系统和内存管理；为什么 OOM 不一定表示机器真的没有任何空闲字节。

这些问题回答清楚，后面的虚拟内存、文件系统和性能分析才有根基。
